\chapter{Calibration, and what is moving}

\textbf{Purpose:} Stop asserting what an instrument can see and measure it, by pushing a wave of
known amplitude through the same code and fading it until the readout loses it. Then ask the one
question every earlier reading threw away by taking a magnitude: not whether anything is large, but
whether anything is \emph{moving}. \textbf{Scope:}
\texttt{src/\allowbreak{}bench/\allowbreak{}bench\_\allowbreak{}sac.\allowbreak{}cu}, the
\texttt{waves} arm.

\section{A null is worth more as a bound}

Every null in this work says the same unsatisfying thing: nothing was found. That is a statement
about the reader as much as about the function, and it is worth nothing until the reader's
sensitivity is known. An instrument that can see a bias of one part in a hundred and one that can
see one part in a hundred thousand both report ``nothing'' on the same data, and only the second
result is interesting.

So the harness is made to carry a plane wave of stated amplitude. A matrix is built that is
binomial by construction --- the same kernel shape, the same shared tally, the same atomics and the
same denominator as the real arm, with the compression replaced by randomness --- and into it a
bias of \(\pm 2^{-(k+1)}\) is driven, positive or negative according to the sign of a wave at each
cell's residue. The bias is built without arithmetic on probabilities: a fair bit or-ed with a rare
bit at \(2^{-k}\) lands on \(0.5 + 2^{-(k+1)}\) exactly, and and-ed with its complement lands on
\(0.5 - 2^{-(k+1)}\).

Fading \(k\) until the readout loses the wave converts ``nothing found'' into ``nothing above this
amplitude''.

\begin{center}
\begin{tabular}{lrl}
\toprule
readout & \(5\sigma\) floor & matched to \\
\midrule
chi-square over all cells & \(1.7 \times 10^{-4}\) & nothing in particular \\
residue fold & \(3.05 \times 10^{-5}\) & residue-structured shapes \\
\bottomrule
\end{tabular}
\end{center}

The fold is about six times sharper than the chi-square, and for a reason that is now measured
and no longer argued: the injected wave is residue-structured and the fold is the matched filter for
it. That ratio is the value of knowing what shape to look for, stated as a number.

\textbf{What this excludes.} No residue-structured wave with a per-cell bias above
\(3.05 \times 10^{-5}\) survives past round twenty-three. The published avalanche dip of
\(0.00272\) is eighty-nine times that floor.

\section{The instrument proves itself}

A calibration is only as good as the evidence that the readout recovered what went in. Injected at
frequency three across the thirty-two residues, the field's autocorrelation comes back as a cosine
of period \(10.7\), and \(32/3 = 10.67\).

\begin{center}
\begin{tabular}{lrrrrrrrr}
\toprule
lag & 1 & 2 & 3 & 4 & 5 & 6 & 7 & 8 \\
\midrule
along output & 45.29 & 16.63 & \(-9.47\) & \(-36.85\) & \(-62.98\) & \(-52.68\) & \(-26.91\) & \(-0.78\) \\
along input & 44.52 & 18.06 & \(-9.05\) & \(-36.51\) & \(-62.95\) & \(-52.56\) & \(-26.84\) & \(0.87\) \\
\bottomrule
\end{tabular}
\end{center}

The two axes agree to three digits, and that agreement is the reflection symmetry of the residue coordinate
\((i - j) \bmod 32\) reading consistently in both directions, and not a separate finding.

\section{The field is white}

A designed function should leave a residual with no length scale at all. A field that is
\emph{smooth} --- correlated between neighboring cells --- would be structure that survived the
mixing, and it is the kind of thing an eye reports on noise readily enough that it has to be
measured and not looked at.

\begin{center}
\begin{tabular}{lrl}
\toprule
field & largest \(|\text{correlation}|\) & verdict \\
\midrule
harness alone & 2.15 & white \\
SHA-256 round 32 & 2.40 & white \\
SHA-256 round 64 & 2.06 & white \\
\bottomrule
\end{tabular}
\end{center}

Sixteen lags per field, eight on each axis, so the largest of them peaks near
\(\sqrt{2 \ln 16} = 2.36\) under a white null whatever the data does. All three sit on it. The
field carries no length scale in either direction, in SHA-256 and in a matrix that is binomial by
construction alike.

\section{What the probe contributes}

The dependency matrix is never the function alone. It is the function convolved with the probe
geometry, and the residue coordinate exists only because input and output bits were indexed by
thirty-two bit words. The rotations act on that same lattice. So the ridge at residue twenty-six is
an interference between the probe's arrangement and the function's, and not a property of
either one, and that bounds what any reading of this shape can be taken to mean.

This is not a defect to be removed. It is the reason the fold works at all: a probe aligned to the
function's own lattice renders a thinly spread structure visible. It is stated here because
an instrument that is part of its own measurement should say so.

\section{What is moving, seen or not}

Every reading up to here looks at one round at a time and asks whether anything in it stands above
the noise. A structure too faint to clear that floor in any single round is invisible to all of
them however many rounds are looked at, because the rounds are compared one at a time and then
discarded.

But the rounds are not independent samples of an unknown thing. They are a deterministic sequence,
so the structure at round \(r\) is related to the structure at \(r+1\) by the round function, and
that relation is carried in the \emph{phase} of the residue spectrum --- which every earlier
reading threw away by taking a magnitude. A structure standing still holds its phase. One drifting
at \(v\) residues per round advances it by \(2\pi f v / 32\) each round. Noise does neither.

Summing the complex spectrum across rounds against a velocity hypothesis therefore adds a real
structure as \(R\) and noise as \(\sqrt{R}\): a gain of \(\sqrt{42} = 6.48\), available only to
whatever is actually moving at the velocity being tried.

\textbf{Each round takes its own seed}, and the separate seeds are what keep this from repeating the
pooling mistake of the previous chapter. The structure is deterministic and stays coherent whatever
the seed; the sampling noise is made independent and decoheres. One shared seed would keep the
noise coherent too and the gain would be nothing.

\subsection{The coordinate the readout lives in}

The scan is run in \(\text{drift} = f v\) and not in velocity. The phase a round advances is
\(2 \pi (f v) r / 32\), so \(fv\) is the natural coordinate and \(v\) is a degenerate one: every
frequency's cone passes through \(v = 0\) together, and away from it the same step in \(v\) is a
step of \(f\) times as much phase.

A uniform grid in \(v\) is therefore correct at one frequency and wrong at all the others. The
first pass used one, and at \(f = 15\) its step was \(1.875\) in drift where forty-two rounds
resolve \(0.195\) --- ten times too coarse to be safe. That mattered, because the first pass's
loudest SHA-256 cell sat at exactly \(f = 15\). Scanning uniformly in drift is the same map with
the degenerate point resolved, and it is one grid that is correct at every frequency at once.

\subsection{The result}

\begin{center}
\begin{tabular}{lrrrr}
\toprule
 & harness alone & walking wave & SHA-256 & null peak \\
\midrule
magnitude & 4.02 & \(\mathbf{131.24}\) & 3.98 & 4.66 \\
phase only & 4.04 & 9.15 & 4.15 & 4.66 \\
best window & 4.07 & 94.58 & 4.37 & 4.66 \\
standing still & 2.19 & 3.50 & 2.25 & --- \\
\bottomrule
\end{tabular}
\end{center}

\textbf{The positive control is the result worth keeping.} The wave was injected at frequency three
walking one residue per round, at an amplitude the ladder had already shown sits below the standing
floor. Its standing-still reading is \(3.50\), \emph{under} the null peak: it genuinely cannot be
found by any single-round instrument at any sample size those instruments could reach. The scan
returned it at \(131.24\), at exactly frequency three and velocity \(-1.000\). That is a factor of
thirty-seven over the standing reading, and all of it came from the motion.

Against an instrument validated that way, nothing in SHA-256's dependency matrix past round
twenty-three is moving, at roughly six and a half times the sensitivity of any standing reading:
an effective floor near \(4.7 \times 10^{-6}\).

\subsection{Two corrections that changed the answer}

Neither is incidental, and both are the same error class this work keeps meeting.

\textbf{Four maxima are not one maximum.} The scan reports four channels and each is itself a
maximum over a grid. Comparing any one of them to the peak of a single channel is the max-of-\(N\)
error one level up. The window scan alone searches \(42{,}570\) cells because sliding the start and
the width multiplies the grid and does not add to it; the union is \(52{,}299\), and the peak
that has to be cleared is \(\sqrt{2 \ln 52299} = 4.66\), not the \(4.12\) of a single channel.
Read against \(4.12\), SHA-256's \(4.37\) would have been announced as a detection. Read correctly
it is flat, and so is the harness at \(4.07\).

\textbf{The phase channel is ceiling-limited.} Summing \(R\) unit vectors cannot exceed \(R\), so
that statistic is bounded above by \(\sqrt{2R}\) whatever the data does --- \(9.17\) at forty-two
rounds. The injected wave reads \(9.15\), which is saturation and not a measurement. Against a
null peak of \(4.66\) the channel has a factor of two of usable range, so discarding amplitude is a
good idea that this baseline cannot afford. It would need many more rounds than SHA-256 has.

